

\section{Lower Bounds}
\label{sec:lower_bounds}

In this section, we present our main conditional hardness results. In Section~\ref{sec:lb_overview}, we outline our approach. In Section~\ref{sec:improved_lower_bounds_effres}, we present our lower bounds for the problem of estimating effective resistances for all pairs of nodes (case where $S = V \times V$), which we call the \emph{all pairs effective resistance estimation problem}. In Section~\ref{sec:improved_lower_bounds_spectral_sums} we show that our techniques also yield lower bounds for other spectral sum estimation problems.
\subsection{Our Approach}\label{sec:lb_overview}
\paragraph{Approaches of Prior Work}
The approach of \cite{musco2017spectrum} begins with the fact that $G$ has a triangle if and only if $\trace{\A_G^3}/6 \geq 1$. They use the fact that various spectral sums $\mathcal{S}_f$ of the SDD matrix $\I - \delta \A_G$ (for $\delta$ sufficiently small) can be expressed as a power series $\mathcal{S}_f(\I - \delta \A_G) = \sum_{k=0}^\infty c_k \delta^k \tr(\A_G^k)$.
The first two terms of this series can be computed directly. So given $Y \approx_\epsilon \mathcal{S}_f(\I - \delta \A_G)$, one can estimate $\trace{\A_G^3}$, where the estimation error is controlled by the magnitude of the first two terms of the series and the tail error due to truncating at the third term. By bounding this estimation error, \cite{musco2017spectrum} shows that, for appropriate choices of $\delta$, $Y \approx_\epsilon \mathcal{S}_f(\I - \delta \A_G)$ yields an additive 1/2 approximation to $\trace{\A_G^3}$, which is sufficient for triangle detection. They also reduce the problem of estimating the spectral sum $\trace{\B^{-1}}$ for an SDD matrix $\B$ to the all pairs effective resistance estimation problem. 

\paragraph{Our Approach} We use three key techniques to better bound the estimation errors incurred in the power-series-inspired approach of \cite{musco2017spectrum}. This yields faster reductions and better lower bounds for effective resistance estimation. Rather than obtaining effective resistance lower bounds by reducing the problem of computing $\trace{\A_G^3}/6$ to computing the \emph{trace} of an SDD matrix as in \cite{musco2017spectrum}, we use a reduction that closely leverages the structure of effective resistances and properties of $\A_G^2$. 

We begin with the fact that for $\alpha>0$ sufficiently small, 
\begin{align}\label{eq:series2}
    \left(\I - \frac{\alpha}{n}\A_G\right)^{-1} = \sum_{k=0}^\infty \frac{\alpha^k}{n^k} \A_G^k. 
\end{align}
Since $\A_G$ is known, given access to $\deltaij^\top (\I - \frac{\alpha}{n}\A_G)^{-1} \deltaij$, we can estimate the entries of $\A_G^2$, where the estimation error is controlled by $\alpha$ and the tail error of truncating \eqref{eq:series2} at the third term. By bounding this estimation error, for appropriate choice of $\alpha$, we can obtain additive 1/2 approximations to all entries of $\A_G^2$, which is sufficient to identify all paths of length 2. We can then detect a triangle by simply scanning for an edge $\{u, v\}$ such that $u$ and $v$ are connected by a path of length 2. Estimating the entries of $\A_G^2$ leads to lower estimation error than estimating $\trace{\A_G^3}$ as in \cite{musco2017spectrum}). 


Second, we use a standard randomized reduction that reduces the triangle detection problem to the triangle detection problem restricted to tripartite graphs. The reduction relies on the fact that a randomly sampled tripartition of the original graph preserves triangles with constant probability. To detect a triangle in a tripartite graph $G = (V_1 \sqcup V_2 \sqcup V_3, E)$, we construct a graph $\Gprime$ by removing all edges $E_{1,2} := \{\{u, v\} \in E : u \in V_1, v \in V_2\}$ between $V_1$ and $V_2$. $G$ has a triangle if and only if there is an edge $\curs{u,v} \in E_{1,2}$ and a path of length 2 between $u$ and $v$ in $\Gprime$. Crucially, we can show that the third term $\As_{\Gprime}^3$ does not contribute to the tail error when estimating the $\{u, v\}$-th entry of $\pars{\A_{\Gprime}^2}$ using \eqref{eq:series2}. 


Third, to lower the spectral norm of $\A_{\Gprime}$ (and consequently better bound the convergence of the power series \eqref{eq:series2}), we use a symmetric random signing of $\A_{\Gprime}$ defined below.
\begin{definition}[Symmetric Random Signing]
    Given a symmetric matrix $\A \in \R^{n \times n}$, its \emph{symmetric random signing} $\As$ is the random matrix with $\As_{i, j} := \rad_{i,j} \A_{i,j}$, where $\rad_{i,j}$ are independent Rademacher random variables (i.e. $\rad_{i,j} = \pm1$, each with probability 1/2) that satisfy $\rad_{i,j} = \rad_{j,i}$.
\end{definition}
We show that with constant probability, this random signing preserves whether the entries of $\A_{\Gprime}^2$ are non-zero, allowing us to detect if $G$ has a triangle even if we replace $\A_G$ in \eqref{eq:series2} with $\As_{\Gprime}$ instead. This is beneficial, as matrix Chernoff guarantees $\Norm{\As_{\Gprime}}_2 = \Otilde{\sqrt{n}}$ whp. whereas $\Norm{\A_{\Gprime}}_2$ may be as large as $n$. This means that by replacing $\A_G$ with $\As_{\Gprime}$, the norm of each term in \eqref{eq:series2} decreases from $O(\alpha^k)$ to $O(\alpha^k n^{-k/2})$. So the tail error of truncating the power series is smaller. 

To compute entries of $\As_{\Gprime}^2$ efficiently using effective resistance estimates on expanders, we first show that we can use all pairs effective resistance estimates on expanders to estimate $\deltaij^T \M^{-1} \deltaij$ for all $i,j \in [n]$, where $\M = (\I - \Q)$ is an SDD matrix with $\rho(\Q) \leq 1/3$. Then, by choosing $\M = \I - \frac{\alpha}{n} \As_{\Gprime}$ as in \eqref{eq:series2} for an appropriate constant $\alpha$, we can estimate $\As_{\Gprime}^2$ from estimates of $\deltaij^T \M^{-1} \deltaij$. This yields our lower bound on the all pairs effective resistance estimation problem.

Additionally, we show that with constant probability, the random signing preserves the property that $\trace{\A_G^3}$ is non-zero. We leverage this aspect of the random signing to obtain improved randomized conditional lower bounds for various spectral sum estimation problems.  We closely follow the trace estimation approach of \cite{musco2017spectrum}, and again use the smaller spectral radius of $\As_G$ to improve bounds on the power series truncation error.

\subsection{Improved Lower Bounds for Effective Resistance Estimation}\label{sec:improved_lower_bounds_effres}
In this section we provide a series of reductions which yield our main result on lower bounds for the all pairs effective resistance estimation problem for all pairs of nodes (case where $S = V \times V$). First, we formalize a standard randomized reduction from triangle detection in general graphs to triangle detection in tripartite graphs in \Cref{lemma:tripartite_to_normal_tridetect} below.

\begin{restatable}{lemma}{lemmaTriangleToTripartite}\label{lemma:tripartite_to_normal_tridetect}
Given an algorithm which can solve the triangle detection problem on an $n$-node tripartite undirected graph in $O(n^\gamma)$ time, we can produce a randomized algorithm which can solve the triangle detection problem on an arbitrary $n$-node undirected graph $G$ in $\widetilde{O}(n^\gamma)$ time whp.
\end{restatable}
\begin{proof}
    We first sample a tripartite subgraph $H$ of $G$ by assigning each vertex in $G$ to a random tripartition with equal probability 1/3 and deleting edges within each resulting tripartition. First, note that if $G$ has no triangles then $H$ also has no triangles since $H$ is a subgraph of $G$. Second, observe that if $G$ has a triangle $\curs{a,b,c}$, it is also a triangle in $H$ if each vertex ends up in a different tripartition. This occurs with probability at least $(1/3)^3 = 1/27$ which is a constant. Therefore, solving the triangle detection problem on $H$ and returning the same output also successfully solves the triangle detection problem on $G$ with probability at least $1/27$. We can repeat this randomized procedure $\log (n^c)$ times to boost the success probability to at least $1 - n^{-c}$, which is whp. in $n$. This randomized algorithm runs in $\widetilde{O}(n^\gamma)$ time, completing the proof.
\end{proof}


\begin{definition}\label{def:Reff_SDD} In the \emph{all pairs SDD effective resistance estimation problem}, given an SDD matrix $\M$ such that $\diagonal{\M} = \I$, $\adjacency{\M} = \Q$, with $\rho(\Q) \leq 1/3$ and $\epsilon \in (0, 1)$, we must output $\widetilde{r} \in \R^{n^2}$ such that $\widetilde{r}_{a,b} \approx_\epsilon \deltaab^\top \M^{-1}\deltaab$ $\forall a, b \in [n]$. We call $\deltaab^\top \M^{-1}\deltaab$ the SDD effective resistance of $(a, b)$ in $\M$. 
\end{definition}

For brevity, we use $\widetilde{r}(\M)$ to refer to the solution of the SDD effective resistance problem on input $\M$. Our first step is to show that an algorithm for the all pairs effective resistance estimation problem on expanders implies an algorithm for the all pairs SDD effective resistance estimation problem.

\begin{restatable}{lemma}{graphEffresToSDDEffres}\label{lemma:grapheffres_to_sddeffres} 
Given an algorithm to solve the all pairs effective resistance estimation problem on graphs with $\widetilde{\Omega}(1)$-expansion in $\Otilde{n^2\epsilon^{-c}}$ time for some $c > 0$, we can produce an algorithm to solve the all pairs SDD effective resistance estimation problem in $\Otilde{n^2 \epsilon^{-c}}$ time. 
\end{restatable}

To prove Lemma~\ref{lemma:grapheffres_to_sddeffres}, we first prove the lemma for the case where $\Q$ is entrywise non-negative in \Cref{lemma:lowerbound_thm2_lemma1} by constructing an expander $G$ with $n+1$ vertices such that $r_{G}(a, b) = \deltaab^\top \M^{-1}\deltaab$ for all $a, b \in [n]$. Then, in \Cref{lemma:lowerbound_thm2_lemma2}, we extend the reduction to arbitrary $\Q$ by constructing $\Q'$ of size $2n$ so that $\Q'$ is entrywise non-negative and $\widetilde{r}(\I - \Q)$ is a simple linear transform of $\widetilde{r}(\I - \Q')$. 

\begin{lemma}\label{lemma:lowerbound_thm2_lemma1} Given an algorithm that solves the all pairs effective resistance estimation problem on graphs with $\widetilde{\Omega}(1)$-expansion in $\Otilde{n^2\epsilon^{-c}}$ time for some $c > 0$, we can produce an algorithm which takes as input an SDD matrix $\M = \I - \Q$ such that $\Q$ is entrywise non-negative and $\Rho{\Q} \leq \frac{1}{3}$, and solves the all pairs SDD effective resistance estimation problem for $\M$ in $\Otilde{n^2 \epsilon^{-c}}$ time.
\end{lemma}

\begin{proof}
Let $v := (\I-\Q) \one$. Note that $v$ is entrywise non-negative. Consider the matrix 
\begin{align*}
    \mathbf{L} := \begin{pmatrix}
        \I & 0 \\
        0 & \Norm{v}_1
    \end{pmatrix} - \begin{pmatrix}
        \Q & v \\
        v^\top & 0
    \end{pmatrix}. 
\end{align*}
and note that it is the Laplacian matrix. Since $\M$ is a principal submatrix of $\mathbf{L}$, we have by the eigenvalue interlacing theorem \cite{hwang2004cauchy} that $\lambda_1 \pars{\M} \leq \lambda_2 \pars{ \mathbf{L} }$. But since $\rho\pars{\Q} \leq 1/3$, we have that $\lambda_1 \pars{\M} \leq 2/3$ and consequently  $\lambda_2\pars{\mathbf{L}} \geq 2/3.$ Therefore, $\mathbf{L}$ is the Laplacian of an $\widetilde{\Omega}(1)$-expander (see Cheeger's Inequality \Cref{thm:cheeger}).

We claim that for $i,j \in n$, 
\begin{align*}
    \deltaij^\top \Ldagger \deltaij = \deltaij^\top \M^{-1} \deltaij, 
\end{align*}
which is sufficient to prove the lemma. Note that for any $x \in \R^n$, $y \in \R$, and $\alpha \in \R$, 
\begin{align*}
    \mathbf{L} \begin{pmatrix}
        x + \alpha \one \\
        y + \alpha
    \end{pmatrix} = \begin{pmatrix}
        (\I - \Q) x + v^\top x \\
        v^\top x - \Norm{v}_1 y 
    \end{pmatrix}. 
\end{align*}
Let 
\begin{align*}
    \Ldagger \deltaij = \begin{pmatrix}
        z_x \\
        z_y
    \end{pmatrix}. 
\end{align*}
Then, note that 
\begin{align*}
    \mathbf{L} \begin{pmatrix}
        z_x - z_y \one \\
        0
    \end{pmatrix} = \begin{pmatrix}
        \deltaij \\
        0
    \end{pmatrix} \implies (\I - \Q) (z_x - z_y \one) = \deltaij. 
\end{align*}

Consequently, 
\begin{align*}
    \deltaij^\top (\I - \Q)^{-1} \deltaij = \deltaij^\top (z_x - z_y \one) = \deltaij^\top z_x - z_y \deltaij^\top \one = \deltaij^\top z_x = \deltaij^\top \Ldagger \deltaij. 
\end{align*}
\end{proof}

Now, to prove \Cref{lemma:lowerbound_thm2_lemma2}, we first show a useful property of block symmetric matrices in the helper lemma below.
\begin{lemma}\label{lemma:lowerbound_thm2_lemma0} Suppose $\A = \begin{pmatrix}
    \X & \Y \\
    \Y & \X
\end{pmatrix}$ where $\X, \Y \in \mathbb{R}^{n \times n}$. Then, 
\begin{align*}
    \deltaij^\top (\X-\Y) \vec{\delta}_{i,j} = \frac{1}{2} \left[\vec{\delta}_{i,n+i}^\top \A \vec{\delta}_{i,n+i} + \vec{\delta}_{j,n+j}^\top \A \vec{\delta}_{j,n+j}\right] - \vec{\delta}_{i,n+j}^\top \A \vec{\delta}_{i,n+j} + \vec{\delta}_{i,j}^\top \A \vec{\delta}_{i,j}. 
\end{align*}
\end{lemma}
\begin{proof} Note that 
\begin{align*}
    \deltaij^\top (\X-\Y) \vec{\delta}_{i,j} &= (\X_{i,i} + \X_{j,j}) - 2 \X_{i,j} + 2\Y_{i,j} - (\Y_{i,i} + \Y_{j,j}). 
\end{align*}
Meanwhile, 
\begin{align*}
    \frac{1}{2} \vec{\delta}_{i,n+i}^\top \A \vec{\delta}_{i,n+i} &= (\X_{i,i} - \Y_{i,i}), \\
    \frac{1}{2} \vec{\delta}_{j,n+j}^\top \A \vec{\delta}_{j,n+j} &= (\X_{j,j} - \Y_{j,j}), \\
    -\vec{\delta}_{i, n+j}^\top \A \vec{\delta}_{i,n+j} &= -\X_{i,i} + 2 \Y_{i,j} - \X_{j,j}, \\
    \deltaij^\top \A \deltaij &= \X_{i,i} - 2 \X_{i,j} + \X_{j,j}, 
\end{align*}
and adding these four terms together concludes the proof. 
\end{proof}

\begin{lemma}\label{lemma:lowerbound_thm2_lemma2} Suppose we are given an algorithm which takes as input an SDD matrix $\M = \I - \Q$ such that $\Q$ is entrywise non-negative and $\Rho{\Q} \leq \frac{1}{3}$, and solves the all pairs SDD effective resistance estimation problem for $\M$ in $\Otilde{n^2 \epsilon^{-c}}$ time. Then, we can produce an algorithm which takes as input an SDD matrix $\Mp = \I - \Qp$ and $\Rho{\Qp} \leq \frac{1}{3}$, and solves the all pairs SDD effective resistance estimation problem for $\Mp$ in $\Otilde{n^2 \epsilon^{-c}}$ time.
\end{lemma}


\begin{proof} We can decompose $\Qp$ as $\Qp = \Ps - \N$ where $\Ps$ is a matrix which contains only the positive offdiagonal entries of $\Qp$ and $-\N$ is a matrix which contains all the negative offdiagonal entries. Therefore both $\Ps$ and $\N$ themselves are entrywise non-negative. We define
\begin{align*}
    \Q :=  \begin{pmatrix}
        \Ps & \N \\
        \N & \Ps
    \end{pmatrix}.
\end{align*}
 
Note that $\Q$ is also entrywise non-negative.


We also have $\rho (\Q) \leq 1/3$. To see this, assume for the sake of contradiction that $\Q$ has an eigenvalue $\lambda > 1/3$. This means that there must exist some $x \in \R^{2n}$ such that $\Q x = \lambda x$. Let $x = [x_1; x_2]$ where $x_1, x_2 \in \mathbb{R}^n$. The eigenvalue equation then implies that $\Ps x_1 + \N x_2 = \lambda x_1$ and $\N x_1 + \Ps x_2 = \lambda x_2$. Subtracting these equations yields $(\Ps - \N)(x_1 - x_2) = \lambda (x_1 - x_2)$ which means that $\lambda$ is also an eigenvalue of $\Qp$. This is a contradiction since $\rho(\Qp) \leq 1/3$. Therefore, $\rho(\Q) \leq 1/3$.

Now, consider the following block decomposition of $\Q^k$
\begin{align*}
    \Q^k = \begin{pmatrix}
        \X & \Y \\ 
        \Y & \X
    \end{pmatrix}
\end{align*}
where $\X, \Y \in \R^{n \times n}$. We will show by induction that $\Qp^k = (\Ps - \N)^k = \X - \Y$. In the base case, when $k = 1$, this is trivially true. Now, assume that the claim holds for all $k \leq m$ for some $m$. Consider the following block decomposition of $\Q^m$
\begin{align*}
    \Q^m = \begin{pmatrix}
        \W & \Z \\ 
        \Z & \W
    \end{pmatrix}.
\end{align*}
By the inductive hypothesis, we know that  $\Qp^m = (\Ps - \N)^m = \W - \Z$. Now, we have that
\begin{align*}
    \Q^{m+1} = \begin{pmatrix}
        \W & \Z \\ 
        \Z & \W
    \end{pmatrix}  \begin{pmatrix}
        \Ps & \N \\ 
        \N & \Ps
    \end{pmatrix} = \begin{pmatrix}
        \W\Ps + \Z\N & \W\N + \Z\Ps \\ 
        \W\N + \Z\Ps & \W\Ps + \Z\N
    \end{pmatrix},
\end{align*}
and also that 
\begin{align*}
    \Qp^{m+1} = (\W - \Z)(\Ps - \N) = \W\Ps - \Z\Ps - \W\N + \Z\N = (\W\Ps + \Z\N) - (\W\N + \Z\Ps). 
\end{align*}
Hence, the claim follows by induction. By Lemma~\ref{lemma:lowerbound_thm2_lemma0}, it follows that 
\begin{align}\label{eq:lemma:lowerbound_thm2_lemma1_eq1}
    \deltaij^\top \Qp^k \deltaij =  \frac{1}{2} \left[\vec{\delta}_{i,n+i}^\top \Q^k \vec{\delta}_{i,n+i} + \vec{\delta}_{j,n+j}^\top \Q^k \vec{\delta}_{j,n+j}\right] - \vec{\delta}_{i,n+j}^\top \Q^k \vec{\delta}_{i,n+j} + \vec{\delta}_{i,j}^\top \Q^k \vec{\delta}_{i,j}. 
\end{align}

Now we can use the power series expansion of $(\I - \Q)^{-1}$ to say, for any $u, v \in [2n]$,
\begin{align*}
    \vec{\delta}_{u,v}^\top (\I - \Q)^{-1} \vec{\delta}_{u,v} = \sum_{k=0}^\infty \vec{\delta}_{u,v}^\top \Q^k \vec{\delta}_{u,v}. 
\end{align*}
Similarly, for any $i, j \in [n]$, 
\begin{align*}
    \vec{\delta}_{i,j}^\top (\I - \Qp)^{-1} \vec{\delta}_{i,j} = \sum_{k=0}^\infty \vec{\delta}_{i,j}^\top \Qp^k \vec{\delta}_{i,j}. 
\end{align*}
So, by linearity and \eqref{eq:lemma:lowerbound_thm2_lemma1_eq1}, it follows that 
\begin{align*}
    \deltaij^\top (\I - \Qp)^{-1} \deltaij =&  \frac{1}{2} \left[\vec{\delta}_{i,n+i}^\top (\I - \Q)^{-1} \vec{\delta}_{i,n+i} + \vec{\delta}_{j,n+j}^\top (\I - \Q)^{-1} \vec{\delta}_{j,n+j}\right]\\ 
    &-\vec{\delta}_{i,n+j}^\top (\I - \Q)^{-1} \vec{\delta}_{i,n+j} + \vec{\delta}_{i,j}^\top (\I - \Q)^{-1} \vec{\delta}_{i,j},
\end{align*}
and this completes the proof.

\end{proof}
\Cref{lemma:grapheffres_to_sddeffres} then follows directly from \Cref{lemma:lowerbound_thm2_lemma1} and \Cref{lemma:lowerbound_thm2_lemma2}. We now turn our attention to reducing the triangle detection problem to the all pairs SDD effective resistance estimation problem. As discussed, a key aspect of our approach is to work with the random signing $\widebar{\A}_G$ of $\A_G$. Lemma~\ref{lemma:signing_2paths} shows that to determine whether $(\A_G^2)_{i,j} > 0$ with constant probability, it suffices to determine whether $(\widebar{\A}_G^2)_{i,j} > 0$. 


\begin{restatable}{lemma}{lemmaSigningTwoPaths}\label{lemma:signing_2paths} For $i \neq j$, if $(\A_{G}^2)_{i,j} = 0$, then $(\As_{G}^2)_{i,j} = 0$; if $(\A_G^2)_{i,j} > 0$, $\mathbb{P}\left[\lvert (\As_G^2)_{i,j}\rvert > 1\right] \geq 1/2$. 
\end{restatable}
\begin{proof} 
The idea of the proof is that if $\{a, b\}, \{b, c\}$ exist in $G$, either $\xi_{a,c} = 1$ or $\xi_{a, c} = -1$ results in $(\As_{G}^2)_{a,c} > 0$. Note that 
\begin{align*}
    \left(\As^2\right)_{i,j} = \sum_{k=1}^n \As_{i,k} \As_{k,j}
\end{align*}

If $\A^2_{i,j} = 0$, then $G$ has no path of length exactly two between $i$ and $j$, so each term $\A_{i,k} \A_{k,j}$ in the summation above must be zero, and hence $\left(\As^2\right)_{i,j} = 0$, completing the first part of the lemma.

Since $\left(\Abar^2\right)_{i,j}$ is only supported on the integers, to prove the second statement, it suffices to show that when $\left(\A^2\right)_{i,j} > 0$, $\mathbb{P}\left[\left\lvert \left(\As^2\right)_{i,j}\right\rvert = 0\right]$ is no larger than 1/2. To see this, note that if $\A^2_{i,j} \neq 0$, then then $G$ has at least one path of length exactly 2 between $i$ and $j$. That is, there exists a $k' \in [n]$ such that $\A_{i,k'} \A_{k', j} = 1$. We can write 
\begin{align*}
    \left(\As^2\right)_{i,j} = \sum_{k = 1}^n \xi_{i, k} \xi_{k,j} \A_{i,k} \A_{k,j}
\end{align*}
Because $i \neq j$, for every $k, \ell \in [n]$ each $\xi_{i,k}$ is always independent from any other $\xi_{\ell, j}$ term appearing in the sum. Moreover, if $\xi_{i,k}$ appears in the sum, then $\xi_{k,i}$ never appears in the sum. Therefore, for each $\{\xi_{i, k} \xi_{k,j}\}_{k\in[n]}$ are themselves independently chosen Rademacher random variables, and all terms in the summation are independent. Separating out the $k'$-th entry,
\begin{align*}
    \left(\As^2\right)_{i,j} = \xi_{i, k'} \xi_{k', j} + \sum_{k \neq k'}^n \xi_{i, k} \xi_{k,j} \A_{i,k} \A_{k,j} \overset{D}{=} \xi' + Z, 
\end{align*}
where $\xi'$ is a Rademacher random variable and $Z := \sum_{k \neq k'}^n \xi_{i, k} \xi_{k,j} \A_{i,k} \A_{k,j}$ is independent from $\xi'$. For any value of $Z$, $\mathbb{P}[\xi' = -Z] \leq \frac{1}{2}$. Therefore, 
\begin{align*}
    \mathbb{P}\left[\left(\As^2\right)_{i,j} = 0\right] \leq \frac{1}{2}, 
\end{align*}
completing the second part of the result. 
\end{proof}

Matrix Chernoff (\Cref{lemma:concentration}) ensures whp. $\rho({\As}_G) = \Otilde{\sqrt{n}}$, while $\rho({\A}_G)$ could be as large as $n$ \cite{tropp2012user}. So, estimating entries of ${\As}_G$ leads to lower power series tail error. 

\begin{restatable}{lemma}{lemmaConcentration}\label{lemma:concentration} Let $G = (V, E)$ be an undirected unweighted graph on $n$ nodes. Let $\As_G$ denote a symmetric random signing of $\A_G$. With high probability, $\rho(\As) \leq \Otilde{\sqrt{n}}$.  
\end{restatable}
\begin{proof} Let $\sigma^2 := \Norm{\sum_{\{u,v\} \in E} \left(\mathbf{E}_{u,v}\right)^2 }_2$ where $\mathbf{E}_{u,v}$ is the adjacency matrix of a graph with only a single edge between $u$ and $v$. Note that entries of $\mathbf{E}_{u,v}^2$ indicates paths of length two in this graph. Therefore, this is a diagonal matrix that satisfies $\left(\mathbf{E}_{u,v}^2\right)_{i,i} = 1$ if and only if $i \in \{u, v\}$. Consequently, 
\begin{align*}
    \sigma^2 = \Norm{\sum_{\{u,v\} \in E} \left(\mathbf{E}_{u,v}\right)^2}_2 = \Norm{\D}_2 \leq \dmax \leq n. 
\end{align*}
We can now write $\As = \sum_{\{u,v\} \in E} \rad_{u,v} \mathbf{E}_{u,v}$ and apply the Matrix Rademacher concentration result (Theorem 1.2) from \cite{tropp2012user}, to get that for any constant $c > 1$, 
\begin{align*}
    \Prob{\lambda_{n}\left( \Abar \right) \geq \sqrt{\dmax} \log(c n)} \leq n \Exp{\frac{-c \dmax \log(n)}{2\dmax}} = n { \frac{1}{n^c} } = n^{-c+1}. 
\end{align*}
\end{proof}

Finally, we use the power series approach in Section~\ref{sec:lb_overview} to obtain \Cref{theorem:sddeffres_hard_as_tridetect}.
\begin{restatable}{theorem}{thmSDDeffresToTripartiteTriangle} \label{theorem:sddeffres_hard_as_tridetect} Given an algorithm which solves the all pairs SDD effective resistance estimation problem in $\Otilde{n^2 \epsilon^{-c}}$ time, we can produce a randomized algorithm that solves the triangle detection problem in $\Otilde{n^{2(1+c)}}$ time whp. 
\end{restatable}
\begin{proof} As $G$ is tripartite, let $V = V_1 \sqcup V_2 \sqcup V_3$ be the partition of $G$ such that no edge has both endpoints in $V_i$ for some $i \in [3]$. Let $E_{1,2} := \{\{u, v\} \in E : u \in V_1, v \in V_2\}$, and let $\Gprime := (V, E\setminus{E_{1,2}})$. Let $\A = \A_{\Gprime}$ denote the adjacency matrix of $\Gprime$ and let $\As$ be a random signing of $\A$.

Suppose that $G$ has a triangle. Then, there exists a pair of vertices $u \in V_1, v \in V_2$ such that $\{u, v\} \in E_{12}$ and $\Gprime$ contains a path of length two between $u$ and $v$. Furthermore, observe that because there are no edges between $V_1$ and $V_2$ in $\Gprime$, $\Gprime$ has no paths of length three between $V_1$ and $V_2$. Consequently, in order to find a triangle in $G$, it suffices to check, for each $i \in V_1, j \in V_2$ with $\{i, j\} \in E_{1,2}$, whether there exists a path of length two between node $i$ and node $j$ in $\Gprime$. In other words, we need to check if there exists some $\curs{i, j} \in E_{1,2}$ such that $\A^2_{i,j} > 0$. By Lemma~\ref{lemma:signing_2paths}, we can instead check if $\abs{\As^2_{i,j}} > 0$ and we would still correctly detect a triangle with probability at least 1/2. Note that this check requires requires only $O(\nnz(\A))$ additional time, since $\Abs{E_{1,2}} < \nnz(\A)$. 


So, our goal now is to compute an accurate enough estimate of $\abs{\As_{i,j}}^2$ for all $\curs{i,j} \in E_{1,2}$ given the effective resistance estimate $ \rtildeij$. To this end, let $\N = \left(\I - \frac{\alpha}{n}\Abar \right)^{-1}$ for some $\alpha < \frac{1}{3}$. Note that the max row-sum of $\Abar$ is 1/3, so the inverse exists. Lemma~\ref{lemma:concentration} guarantees that with high probability, $\Rho{\Abar} \leq \Otilde{\sqrt{\dmax}} = \Otilde{\sqrt{n}}$. We condition on this event in the remainder of the proof. Consequently, we can express $\N$ as a power series,
\begin{align*}
    \N = \sum_{k=0}^\infty \left(\frac{\alpha}{n}\right)^k \Abar^k. 
\end{align*}
Now, let $\Ntilde$ denote the truncation of $\N$ at the third term in this power series. That is, $\Ntilde = \I + \frac{\alpha}{n} \Abar + \frac{\alpha^2}{n^2} \Abar^2$. Therefore, we have for all $\curs{i,j} \in E_{1, 2}$, 
\begin{align*}
    \As_{i,j}^2 = \frac{n^2}{\alpha^2}\Ntilde_{i, j}.
\end{align*}
Noticing that $\N_{i,j} = \frac{\N_{i,i} + \N_{j,j} - r_{i,j}}{2} $ 
 motivates us to define our estimate of $\As^2_{i,j}$ that we denote $P_{i,j}$ as follows
\begin{align*}
    P_{i,j} := \frac{n^2}{\alpha^2} \left[ \frac{\Ntilde_{i,i} + \Ntilde_{j,j} - \rtildeij}{2} \right].
\end{align*}
Now, observe that $\Ntilde_{i,i} = 1 + \frac{\alpha^2}{n^2}\As^2_{i,i}$, and $\As^2_{i,i}$ is simply the degree of vertex $i$ in $\Gprime$. Note that the random signing does not affect the fact that the diagonal entries of the square of the adjacency matrix are the degrees. Therefore, $\Ntilde_{i,i}$ can also be computed for all $i$ in only $O(\nnz(\A))$ time. The additive error between our estimate $P_{i,j}$ and $\abs{\As}^2_{i,j}$ takes the form
\begin{align*}
    \abs{P_{i,j} - \As^2_{i,j}} = \abs{\frac{n^2}{\alpha^2}\sqrs{\Ntilde_{i,j} - \frac{\Ntilde_{i,i} + \Ntilde_{j,j} - \rtildeij}{2}}}.
\end{align*}
By triangle inequality and plugging in the definition of $r_{i,j}$, we break up the error into four pieces
\begin{align}\label{eq:error_in_four_pieces}
    \abs{P_{i,j} - \As^2_{i,j}} \leq \frac{n^2}{\alpha^2} \sqrs{\abs{\Ntilde_{i, j} - \N_{i, j}} + \abs{\Ntilde_{i, i} - \N_{i, i}}/2 + \abs{\Ntilde_{j, j} - \N_{j, j}}/2 +\abs{\rtildeij - r_{i, j}}/2 }.
\end{align}
We now bound each term separately. Consider any $i \in V_1$. Since, $\Gprime$ contains no triangle containing $i$, $\left(\Abar^3\right)_{i,i} = 0$. So, we have that 
\begin{align*}
    \left\lvert \N_{i,i} - \Ntilde_{i,i}\right\rvert &= \left\lvert \sum_{k=4}^\infty \frac{\alpha^k}{n^k} \left(\Abar^k\right)_{i,i}\right\rvert \leq \Norm{ \sum_{k=4}^\infty \frac{\alpha^k}{n^k} \Abar^k }_2 \\
    &\leq \sum_{k=4}^\infty \widetilde{O}\pars{\frac{\alpha}{\sqrt{n}}}^k = \frac{\Otilde{\frac{\alpha}{\sqrt{n}}}^4}{1 - \Otilde{\frac{\alpha}{\sqrt{n}}}}
\end{align*}
Similarly, for any $i \in V_1$ and $j \in V_2$, note that $\Abar^3_{i,j} = 0$ because $\Gprime$ has no edges between $V_1$ and $V_2$, and there are no edges within each tripartition $V_i$. Consequently, by a similar argument as above, 
\begin{align*}
    \left\lvert \N_{i,j} - \Ntilde_{i,j}\right\rvert &= \sum_{k=4}^\infty (\Abar^4)_{i,j} = \frac{1}{2}\sum_{k=4}^\infty e_i^\top (\Abar^k) e_i  +  e_j^\top (\Abar^k) e_j  - \vec{\delta}_{i,j}^\top (\Abar^k) \vec{\delta}_{i,j} \\
    &= \frac{1}{2} \left( e_i^\top e_i + e_j^\top e_j + \vec{\delta}_{i,j}^\top \vec{\delta}_{i,j}\right) \Norm{ \sum_{k=4}^\infty \frac{\alpha^k}{n^k} \Abar^k }_2 \\
    &\leq 2 \sum_{k=4}^\infty \widetilde{O}\pars{\frac{\alpha}{\sqrt{n}}}^k = \frac{2 \Otilde{\frac{\alpha}{\sqrt{n}}}^4}{1 - \Otilde{\frac{\alpha}{\sqrt{n}}}}. 
\end{align*} 

Finally, consider the magnitude of the approximation error between $\rtildeij$ and $r_{i,j}$. We have
\begin{align*}
    \abs{r_{i,j} - \rtildeij} \leq \epsilon \abs{r_{i,j}}.
\end{align*}
Note that $\abs{r_{i,j}} = \abs{\deltaij^\top \N \deltaij} \leq \norm{\deltaij}^2 \norm{\N} \leq 2 \norm{\pars{\I - \frac{\alpha}{n}\As}^{-1}} \leq \frac{2}{1-\frac{\alpha}{n} \Norm{\Abar}_2 } \leq \frac{2}{1 - \Otilde{\frac{\alpha}{\sqrt{n}}}}. $
Plugging these estimates of the errors into \eqref{eq:error_in_four_pieces}, we get
\begin{align*}
    \abs{P_{i,j} - \As^2_{i,j}} \leq \frac{n^2}{\alpha^2} \sqrs{\frac{\widetilde{O}\pars{\frac{\alpha}{\sqrt{n}}}^4 + \epsilon}{1 - \widetilde{O}\pars{\frac{\alpha}{\sqrt{n}}}}}.
\end{align*}

Since $\Abs{\Abar^2_{i,j}} \in \{0, 1\}$, to compute $\Abar^2_{i,j}$, we need approximate it to additive 1/2 error. That is, we require 
\begin{align*}
    \frac{1}{1 - \Otilde{\frac{\alpha}{\sqrt{n}}}} \left[{\Otilde{\alpha^2}} + \epsilon \frac{n^2}{\alpha^2} \right] &< \frac{1}{2},
\end{align*}
or equivalently, 
\begin{align*}
    \epsilon &< \frac{{\alpha^2 - \Otilde{\frac{\alpha^3}{\sqrt{n}}}}}{2n^2} - \frac{\Otilde{\alpha^4}}{n^2} = \widetilde{O}\pars{\frac{1}{n^2}},
\end{align*}
where the last step follows from the fact that we can take $\alpha$ to be a sufficiently small constant. Therefore, estimates $\rtildeij$ with $\epsilon = \widetilde{O}(n^{-2})$ for all $\curs{i, j} \in E_{1,2}$ are sufficient to determine if $\Abs{\Abar^2_{i,j}}$ is 0 or 1. As noted earlier, checking this for all edges in $E_{1,2}$ takes $O(\nnz (\A))$ additional time. Therefore, by plugging in $\epsilon = \widetilde{O}(n^{-2})$, we can use an algorithm that solves the all pairs SDD effective resistance estimation problem in $\widetilde{O}(n^2 \epsilon^{-c})$ time to solve the triangle detection problem in $\widetilde{O}(n^2 n^{2c}) = \widetilde{O}(n^{2(1+c)})$ time whp. and this completes the proof.  
\end{proof}

\Cref{theorem:sddeffres_hard_as_tridetect} and \Cref{lemma:grapheffres_to_sddeffres} with $c = 1/2 - \delta$ immediately imply our main result \Cref{thm:reff_lb_main_result}.

\subsection{Improved Lower Bounds for Spectral Sum Estimation} \label{sec:improved_lower_bounds_spectral_sums}
Finally, we discuss our improved lower bounds for various spectral sum estimation problems. Analogous to Lemma~\ref{lemma:signing_2paths}, in the following lemma we show that to determine whether a graph has a triangle (i.e., $\trace{\A_{G}^3} > 0$) with constant probability, it suffices to determine whether $\trace{\As_{G}^3} > 0$.

\begin{restatable}{lemma}{lemmaRandomSigningPreservesTriangles}\label{lemma:random_signing_preserves_triangles}
    If $\trace{\A_G^3} = 0$, then $\trace{\As_G^3} = 0$, and if $\trace{\A_G^3} > 0$ then $\Prob{\abs{\trace{\As_G^3}} > 0} \geq 1/4$. 
\end{restatable}

\begin{proof}
    The central idea of the proof is that if a triangle $\{a, b\}, \{b, c\}, \{c, a\}$ exists in $G$, then amongst the 4 possible configurations of the Rademacher random signing variables $\xi_{a,b}$ and $\xi_{b,c}$, at least one configuration must result in $\abs{\trace{\As_G^3}} > 0$. 

    Again for convenience, let $\As = \As_G$. Denote by $\rad_{ab} = \rad_{ba}$ the Rademacher random variable used to decide the sign of edge $(a,b)$. We can write
    \begin{align}\label{eq:sum_over_all_triangles}
        \trace{\As^3}/6 = \sum_{\text{triangles } \curs{i,j,k} \text{ in } G}  \rad_{ij} \rad_{jk} \rad_{ki} =: T
    \end{align}
    If $\trace{\A^3} > 0$, then $G$ must have at least one triangle. Consider the following cases:
    \begin{enumerate}
        \item $G$ has an odd number of triangles. In this case, since each term in the sum in \eqref{eq:sum_over_all_triangles} is either $+1$ or $-1$, and there is an odd number of terms in the sum, so $\trace{\As^3} > 0$ wp $1$.
        \item $G$ has an even number of triangles. First, define 
        \begin{align*}
            T(a, b) := \sum_{\text{triangles } \curs{a,b,k} \text{ that contain }\curs{a,b}} \rad_{ab} \rad_{bk} \rad_{ak}
        \end{align*}
        We now subdivide this into two cases:
        \begin{enumerate}
            \item There exists an edge $\curs{a,b}$ that is part of an odd number of triangles. We can decompose the sum in \eqref{eq:sum_over_all_triangles} as follows:
            \begin{align*}
                \trace{\As^3}/6 = \underbrace{T(a,b)}_{S_1} + \underbrace{T-T(a,b)}_{S_2}
            \end{align*}
            Suppose there exists some realization of the random variables $\rad_{ij}$ such that $\trace{\As^3}/6 = 0$. Since $S_1$ has odd terms, it must be non-zero. By flipping the sign of $\rad_{ab}$, we can flip the sign of $S_1$, and so $-S_1 + S_2 \neq 0$. Therefore, for every configuration the variables $\rad_{ij}$ that result in a 0 trace, there exists an equally likely configuration that results in $\trace{\As^3}/6 \neq 0$. Therefore, $\Prob{\abs{\trace{\As^3}} > 0} \geq 1/2$.
            \item Every edge in $G$ is a part of an even number of triangles. Let $\curs{a,b,c}$ be a triangle in $G$. In this case, we decompose \eqref{eq:sum_over_all_triangles} as follows:
            \begin{align*}
                \trace{\As^3}/6 &= \underbrace{\rad_{ab} \rad_{bc} \rad_{ac}}_{S_1} +\underbrace{T(a,b) - \rad_{ab} \rad_{bc} \rad_{ac}}_{S_2} + \underbrace{T(b,c) - \rad_{ab} \rad_{bc} \rad_{ac}}_{S_3}  \\ &+\underbrace{T - T(a,b) - T(b,c) + \rad_{ab} \rad_{bc} \rad_{ac}}_{S_4}
            \end{align*}
            Consider all 4 possible values of the pair of random variables $\curs{\rad_{ab}, \rad_{bc}}$. Since each $S_i$ has an odd number of terms, $S_i \neq 0$ for all $i$. We observe that it is not possible for all 4 equally likely configurations of $\curs{\rad_{ab}, \rad_{bc}}$ to result in $S_1 + S_2 + S_3 + S_4 = 0$, so at least one configuration must result in $T \neq 0$. Therefore,  $\Prob{\abs{\trace{\As^3}} > 0} \geq 1/4$.
        \end{enumerate}
    \end{enumerate}
\end{proof}

By following the proof of Theorem 15 from \cite{musco2017spectrum}, and replacing their use of $\A_G$ with a symmetric random signing $\As_G$, we obtain an improved randomized version of their result by leveraging the smaller spectral radius of $\As_G$: 
\begin{restatable}[Improved randomized version of Theorem 15 from \cite{musco2017spectrum}]{theorem}{thmFifteenSigned}\label{thm:thm15_signed}
Let $f:\R^+ \to \R^+$ be a function such that it can be expressed as $f(x) = \sum_{k=0}^\infty c_k (x - 1)^k$ where $\abs{c_k/c_3} \leq h^{k-3}$ for $k > 3$ and $x \in (0,2)$. Given an algorithm which takes as input a graph $G = (V, E)$ on $n$ nodes and, in $O\pars{n^\gamma \epsilon^{-c}}$ time, outputs an estimate $X \approx_{\epsilon_1/9} \mathcal{S}_f(\I - \delta \As_G)$ with $\delta$ and $\epsilon_1$ satisfying 
\begin{align*}
    &\delta = \min\curs{({\sqrt{n} \log (\alpha n)})^{-1}, ({10n^3 h \log (\alpha n)})^{-1}}\\
    &\epsilon_1 = \min\curs{1, \abs{{c_3 \delta^3}/{(c_0 n)}}, \abs{{c_3 \delta}/{(c_2 n^2)}}}
\end{align*} 
for some constant $\alpha>1$, we can produce an algorithm that solves the triangle detection problem in $O\pars{n^2 + n^\gamma \epsilon_1^{-c}}$ time whp.
\end{restatable}
\begin{proof}
For convenience, we write $\A = \A_G$. The proof closely follows the proof of Theorem 15 from \cite{musco2017spectrum}, but we replace $\A$ with its symmetrically random signed version that we denote $\As$. We present a full proof here for completeness.

First, we note that by lemma~\ref{lemma:concentration}, we have that $\norm{\As}_2 \leq \sqrt{n} \log{\alpha n}$ whp. for some constant $\alpha >1$. We define $\Bs = \I - \delta \As$ and consequently $\Bs$ is PSD whp. Now, using the definition of $f$, we have
    \begin{align*}
        \sum_{i=1}^n \sigma_i(\Bs) = \sum_{i=1}^n f(1 - \delta \lambda_i(\As)) = \sum_{i=1}^n \sum_{k=0}^\infty c_k (\delta \lambda_i(\As))^k = \sum_{k=0}^\infty c_k \delta^k \trace{\As^k}.
    \end{align*}
We analyze the tail of this power series. Specifically, we have 
\begin{align}\label{eq:thm15proof_powerseries_tail_bound}
    \abs{\sum_{k=4}^\infty  c_k \delta^k \trace{\As^k}} \leq \abs{c_3} \delta^3 \sum_{k=4}^\infty \abs{\trace{\As^k}} \delta^{k-3} \abs{\frac{c_k}{c_3}}.
\end{align}
Now, we have $\abs{\trace{\As^k}} \leq \norm{\As}^{k-2}_2 \normgen{\As}{F} \leq n^{k/2 + 1}$ whp. and further $n^{k/2 + 1} \leq n^{3\pars{k-3}}$ for all $k>3$. Therefore, using the definition of $\delta$ as in the theorem we get that whp.,
    \begin{align*}
        \abs{\trace{\As^k}} \delta^{k-3} \abs{\frac{c_k}{c_3}} \leq \frac{1}{10^{k-3}}\text{ for all } k>3.
    \end{align*}
Plugging into \Cref{eq:thm15proof_powerseries_tail_bound}, we get 
\begin{align*}
    \abs{\sum_{k=4}^\infty  c_k \delta^k \trace{\As^k}} \leq \frac{\abs{c_3}\delta^3}{9}.
\end{align*}
The rest of the proof is essentially identical to the steps in the proof of Theorem 15 in \cite{musco2017spectrum}, but we reproduce them here for completeness.

Using the simple facts $\trace{\As^0} = n$, ${\trace{\As}} = 0$ and ${\trace{\As^2}} \leq n^2$, we have
\begin{align*}
    c_0\trace{\As^0} + c_1\trace{\As} + c_2 {\trace{\As^2}} \leq \abs{c_3}\delta^3 \pars{c_0 n/(c_3\delta^3) + c_2 n^2/(c_3 \delta)} \leq \frac{\abs{c_3} \delta^3}{\epsilon_1}.
\end{align*}
Given $X$, in $O(\nnz(\As))$ time, we can compute 
\begin{align*}
    X - c_0 n - c_2 \delta^2 {\trace{\As^2}} &= {c_3}\delta^3 \trace{\As^3} \pm \frac{\abs{c_3} \delta^3}{9} \pm \frac{\epsilon_1}{9}\pars{\frac{\abs{c_3}\delta^3}{9} + c_3 \delta^3 \trace{\As^3} +\frac{\abs{c_3}\delta^3}{\epsilon_1}}\\
    &= c_3 \delta^3 \sqrs{\trace{\As^3}\pars{1 \pm \frac{1}{20}} \pm \frac{1}{3}}.
\end{align*}
This is sufficient to detect if $\abs{\trace{\As}} = 0$ or if $\abs{\trace{\As^3}} \geq 1$. The final result then follows by applying Lemma~\ref{lemma:random_signing_preserves_triangles}.   
\end{proof}

We now prove \Cref{thm:our_results_spectralsums_lowerbound} by applying this result to the functions $f$ that define the corresponding spectral sums.

\begin{proof}[Proof of \Cref{thm:our_results_spectralsums_lowerbound}]
    We apply \Cref{thm:thm15_signed} to the specific spectral sums in \Cref{table:spectral_sum_hardness}.
    \paragraph{Schatten 3-norm} We have $f(x) = x^3$. Therefore, $c_k = 0$ for $k > 3$. So we apply Theorem \ref{thm:thm15_signed} with $h = 0$ and hence $\delta = \Otilde{\frac{1}{\sqrt{n}}}$ and $\epsilon_1 = \Otilde{\frac{1}{n^{2.5}}}$.
    
    \paragraph{Schatten p-norm $p \neq 1, 2$} We have $f(x) = x^p$. Using the Taylor series about 1, we have $\frac{c_k}{c_3} \leq p^{k-3}$ for all $k>3$ as well as $\abs{\frac{c_0}{c_3}} = \abs{\frac{1}{p(p-1)(p-2)}} \leq \abs{\frac{1}{2 \min \curs{p, (p-1), (p-2)}}}$ and similarly $\abs{\frac{c_2}{c_3}} \leq \abs{\frac{1}{2 \min \curs{p, (p-1)}}}$. Therefore, with $h=p$, we apply Theorem \ref{thm:thm15_signed} with $\delta = \Otilde{\frac{1}{n^3 p}}$ and $\epsilon_1 = \frac{c_3 \delta^3}{c_0 n} = \Otilde{\frac{\abs{\min \curs{p, (p-1), (p-2)}}}{n^{10}p^3}}$, which gives the result.

    \paragraph{SVD Entropy} We have $f(x) = x \log x$. For $x \in (0, 2)$, using the Taylor Series about 1 we can write $x \log x =\sum_{k=0}^\infty c_k(x-1)^k$ where $c_0 = 1 \log(1) = 0$, $c_1 = \log(1) + 1 = 1$, and $|c_k| = \frac{(k-2)!}{k!} \leq 1$ for $k \geq 2$. So we have $c_k < c_3$ for all $k > 3$, $\frac{c_0}{c_3} = 0$ and $\frac{c_2}{c_3} = \frac{1}{3}$. So with $h=1$,  Applying Theorem \ref{thm:thm15_signed} with $\delta = \Otilde{\frac{1}{n^3}}$ and $\epsilon_1 = \frac{\delta}{3n^2} = \Otilde{\frac{1}{n^5}}$ gives the result.

    \paragraph{Log Determinant} We have $f(x) = \log x$. For $x \in (0, 2)$, using the Taylor Series about 1 we can write $\log x =\sum_{k=0}^\infty c_k(x-1)^k$ where $c_0 = 0$, $\abs{c_i} = 1/i$ for $i \geq 1$. Again we have $c_k < c_3$ for all $k > 3$ and $\frac{c_0}{c_3} = 0$ while $\frac{c_2}{c_3} = \frac{3}{2}$. So with $h=1$,  Applying Theorem \ref{thm:thm15_signed} with $\delta = \Otilde{\frac{1}{n^3}}$ and $\epsilon_1 = \frac{\delta}{3n^2} = \Otilde{\frac{1}{n^5}}$ gives the result.

    \paragraph{Trace of Exponential} We have $f(x) = e^x$. Using the Taylor Series about $1$ we can write $e^x = \sum_{k=0}^\infty \frac{e(x-1)^k}{k!}$. We have $\frac{c_0}{c_3} = 6$, $\frac{c_2}{c_3} = 3$, and $c_k < c_3$ for all $k \geq 3$. So with $h=1$, Applying Theorem \ref{thm:thm15_signed} with $\delta = \Otilde{\frac{1}{n^3}}$ and $\epsilon_1 = \frac{c_3 \delta^3}{c_0 n} = \Otilde{\frac{1}{n^{10}}}$ gives the result.
\end{proof}